\documentclass[12pt,a4paper]{article}

% --- Codificação e língua ---
\usepackage[utf8]{inputenc}
\usepackage[T1]{fontenc}
\usepackage[portuguese]{babel}

% --- Tipografia ---
\usepackage{lmodern}
\usepackage{setspace}
\onehalfspacing

% --- Margens ---
\usepackage[top=2.5cm, bottom=2.5cm, left=2.5cm, right=2.5cm]{geometry}

% --- Tabelas ---
\usepackage{booktabs}
\usepackage{float}

% --- Cores e hiperligações ---
\usepackage[dvipsnames]{xcolor}
\usepackage[colorlinks=true, linkcolor=NavyBlue, citecolor=NavyBlue, urlcolor=NavyBlue]{hyperref}

% --- Matemática ---
\usepackage{amsmath}

% ============================================================
\begin{document}
% ============================================================

\begin{center}
  {\LARGE\bfseries Resultados de Mutation Testing}\\[0.5em]
  {\large Gonçalo Pires --- Dissertação de Mestrado --- Maio de 2026}
\end{center}

\bigskip
\hrule
\bigskip

% ============================================================
\section{O que foi feito}
% ============================================================

Neste trabalho foi aplicada a técnica de \textit{mutation testing} com a ferramenta \textbf{PIT}
(versão~1.15.8) para avaliar a qualidade das suites de testes desenvolvidas.
O processo consistiu em três etapas principais:

\begin{enumerate}
  \item \textbf{Caso controlado} --- execução do PIT sobre a classe \texttt{dataloader.MathOperations},
        combinando várias suites de testes (manuais e geradas pelo EvoSuite) e refinando-as
        iterativamente com base nos mutantes que sobreviviam.
  \item \textbf{Benchmark SF110} --- execução do PIT sobre os 111 projectos Java do benchmark
        SF110, no modo \texttt{both} (testes manuais + EvoSuite combinados), avaliando assim
        a eficácia do pipeline em projectos reais de dimensão variável.
  \item \textbf{Análise dos resultados} --- comparação por projecto, identificação dos melhores
        scores e discussão das limitações encontradas (incompatibilidades de build, timeouts, etc.).
\end{enumerate}

% ============================================================
\section{Resultados em \texttt{MathOperations}}
% ============================================================

O caso controlado permitiu refinar as suites de testes de forma iterativa.
A combinação de todas as suites disponíveis (\texttt{ALL\_COMBINED}, 229~testes) atingiu
\textbf{94,12\%} de \textit{mutation score} num conjunto de 102~mutações.
Após análise dos mutantes sobreviventes e introdução de melhorias cirúrgicas nas suites,
um rerun focado exclusivamente em \texttt{MathOperations} produziu os seguintes resultados:

\begin{table}[H]
  \centering
  \caption{Rerun focado em \texttt{MathOperations} após refinamento das suites.}
  \label{tab:rerun}
  \begin{tabular}{lrrrrrr}
    \toprule
    \textbf{Testes} & \textbf{Killed} & \textbf{Survived} &
    \textbf{No Cov.} & \textbf{Total} & \textbf{Score (\%)} & \textbf{Strength (\%)} \\
    \midrule
    84 & 71 & 0 & 2 & 73 & \textbf{97,26} & \textbf{100,00} \\
    \bottomrule
  \end{tabular}
\end{table}

O score de \textbf{97,26\%} e a \textit{test strength} de \textbf{100\%} significam que todos
os mutantes cobertos pelos testes foram detectados. Os 2 mutantes restantes estão na categoria
\textit{no coverage} --- correspondem a lambdas geradas pelo compilador em expressões
\texttt{orElseThrow} que não são atingíveis por testes unitários convencionais.

% ============================================================
\section{Resultados no Benchmark SF110}
% ============================================================

A campanha principal (\texttt{sf110\_full\_both1}) correu o PIT sobre os 111 projectos do SF110
em modo \texttt{both}, combinando testes manuais e testes gerados pelo EvoSuite.
Dos 111 projectos, \textbf{79 concluíram com sucesso} (71,2\%). Os restantes 32 falharam
maioritariamente por incompatibilidades de classpath ou ausência de testes, o que é independente
da qualidade do pipeline.

A Tabela~\ref{tab:top20} apresenta os \textbf{20 melhores projectos} por \textit{mutation score}.

\begin{table}[H]
  \centering
  \caption{Top~20 projectos SF110 por \textit{mutation score} --- campanha \texttt{sf110\_full\_both1}.}
  \label{tab:top20}
  \footnotesize
  \begin{tabular}{clrrrrrr}
    \toprule
    \textbf{\#} & \textbf{Projecto} & \textbf{Score (\%)} & \textbf{Strength (\%)} &
    \textbf{Killed} & \textbf{Survived} & \textbf{No Cov.} & \textbf{Total} \\
    \midrule
    1  & \texttt{90\_dcparseargs}  & \textbf{90,32} & 90,32  &  84 &   9 &   0 &   93 \\
    2  & \texttt{31\_xisemele}     & \textbf{90,24} & 91,64  & 296 &  27 &   5 &  328 \\
    3  & \texttt{24\_saxpath}      & \textbf{84,09} & 87,68  & 370 &  52 &  18 &  440 \\
    4  & \texttt{26\_jipa}         & \textbf{75,86} & \textbf{100,00} &  22 &   0 &   7 &   29 \\
    5  & \texttt{40\_glengineer}   & \textbf{67,10} & 74,59  & 414 & 141 &  62 &  617 \\
    6  & \texttt{59\_mygrid}       & \textbf{65,24} & 89,84  & 610 &  69 & 256 &  935 \\
    7  & \texttt{55\_lavalamp}     & \textbf{63,97} & 71,49  & 158 &  63 &  26 &  247 \\
    8  & \texttt{25\_jni-inchi}    & \textbf{62,03} & 85,47  & 147 &  25 &  65 &  237 \\
    9  & \texttt{86\_at-robots2-j} & \textbf{61,90} & 68,29  & 450 & 209 &  68 &  728 \\
    10 & \texttt{2\_a4j}           & \textbf{54,35} & 75,00  & 225 &  75 & 114 &  414 \\
    11 & \texttt{21\_geo-google}   & \textbf{54,23} & 58,58  & 314 & 222 &  43 &  579 \\
    12 & \texttt{43\_lilith}       & \textbf{53,48} & 83,25  & 969 & 195 & 648 & 1812 \\
    13 & \texttt{80\_wheelwebtool} & \textbf{51,99} & 88,26  & 248 &  33 & 196 &  477 \\
    14 & \texttt{39\_diffi}        & \textbf{47,95} & 61,76  & 105 &  65 &  49 &  219 \\
    15 & \texttt{77\_io-project}   & \textbf{44,70} & \textbf{96,72}  &  59 &   2 &  71 &  132 \\
    16 & \texttt{14\_omjstate}     & \textbf{40,91} & 72,00  &  36 &  14 &  38 &   88 \\
    17 & \texttt{13\_jdbacl}       & \textbf{40,82} & 83,02  & 396 &  81 & 493 &  970 \\
    18 & \texttt{101\_netweaver}   & \textbf{39,60} & \textbf{95,06}  &1655 &  86 &2438 & 4179 \\
    19 & \texttt{61\_noen}         & \textbf{38,04} & 60,83  & 264 & 170 & 260 &  694 \\
    20 & \texttt{18\_jsecurity}    & \textbf{37,55} & 75,11  & 175 &  58 & 233 &  466 \\
    \bottomrule
  \end{tabular}
\end{table}

Alguns aspectos a destacar:

\begin{itemize}
  \item \texttt{90\_dcparseargs} e \texttt{31\_xisemele} ultrapassaram \textbf{90\%} de
        \textit{mutation score}, o que é excepcional para projectos reais de código aberto.

  \item \texttt{24\_saxpath} atingiu \textbf{84,09\%} com 440 mutações --- um projecto de
        dimensão média com cobertura muito elevada.

  \item \texttt{26\_jipa} e \texttt{53\_shp2kml} alcançaram \textbf{100\% de test strength}:
        todos os mutantes cobertos pelos testes foram mortos, sem nenhum sobrevivente.

  \item \texttt{77\_io-project} teve \textbf{96,72\% de test strength} e
        \texttt{101\_netweaver} atingiu \textbf{95,06\% de test strength} --- este último com
        o maior número de mutações relevantes da campanha (4\,179 mutações).

  \item Projectos como \texttt{43\_lilith} (53,48\%, 1\,812 mutações) e
        \texttt{59\_mygrid} (65,24\%, 935 mutações) demonstram que o pipeline mantém
        boa eficácia mesmo em projectos de maior dimensão.

  \item Os primeiros 13 projectos da tabela superaram o limiar de \textbf{50\%} de
        \textit{mutation score}, e todos os 20 apresentados superam \textbf{37\%}.
\end{itemize}

% ============================================================
\section{Conclusões}
% ============================================================

Os resultados obtidos demonstram que o pipeline de \textit{mutation testing} desenvolvido é
eficaz tanto em contexto controlado como em projectos reais:

\begin{itemize}
  \item No caso \texttt{MathOperations}, o refinamento iterativo elevou o score de
        \textbf{94,12\%} para \textbf{97,26\%}, com \textbf{test strength de 100\%}.
  \item No SF110, a melhor campanha (\texttt{sf110\_full\_both1}) completou 79 de 111 projectos,
        com vários projectos acima de 60\% e dois acima de \textbf{90\%}.
  \item A combinação de testes manuais com testes gerados automaticamente pelo EvoSuite
        (modo \texttt{both}) revelou-se consistentemente superior ao uso exclusivo de cada
        fonte individualmente.
\end{itemize}

\end{document}
